% Source: https://artofproblemsolving.com/community/c3h3156870_2014_lmt_spring_guts_round_p16p30__lexington_mathematical_tournament
% Stable identity: aops-topic-3156870-q24
Reduce entries modulo 3. There are five residue-zero entries, six residue-one entries, and five residue-two entries. Adjacent residues $ u,v $ are permitted exactly when $(u+v)\bmod3\ne0 $. A row-state recurrence counts all residue patterns with these totals, giving 32 patterns:
\begin{verbatim}
from itertools import product
from collections import defaultdict
from math import factorial
rows=[r for r in product(range(3),repeat=4)
      if all((r[i]+r[i+1])%3 for i in range(3))]
cnt={r:tuple(r.count(i) for i in range(3)) for r in rows}
D={(r,*cnt[r]):1 for r in rows}
for step in range(3):
    E=defaultdict(int)
    for (r,a,b,c),ways in D.items():
        for s in rows:
            if all((x+y)%3 for x,y in zip(r,s)):
                x,y,z=cnt[s]
                if a+x<=5 and b+y<=6 and c+z<=5:E[(s,a+x,b+y,c+z)]+=ways
    D=E
N=sum(w for (r,a,b,c),w in D.items() if (a,b,c)==(5,6,5))
assert N==32
assert N*factorial(5)**2*factorial(6)==331776000
\end{verbatim}
Each pattern allows $ 5!6!5!$ independent assignments of the distinct integers of each residue. Therefore the answer is $\boxed{331776000}$. 
\begin{asy}
usepackage("amsmath");
size(7cm);
for(int k=0;k<=4;++k){draw((k,0)--(k,4));draw((0,k)--(4,k));}draw((.5,.5)--(1.5,.5),blue,Arrows);draw((.5,.5)--(.5,1.5),blue,Arrows);label("Adjacency: shared sides only",(2,-.7));
\end{asy}
